% ============================================================
%  03 · Diskrete Zufallsprozesse
%      Quelle: 03_diskrete_zufallsprozesse.pdf
% ============================================================

\bereich[cDiskret]{Zeitdiskreter Zufallsprozess — Definition}
\begin{formelreihe}
  \formelblock{Abtastung mit Periode $T$}{X_\eta[k]=X(\eta,kT),\quad k\in\mathbb{Z}\\[0.15em]
    \text{Musterfunktion: }X_{\eta_i}[k]=x_i[k]} &
  \infoblock{\textbf{Wichtig:} zeitdiskret $\neq$ wertdiskret!\\
    Werte sind weiterhin kontinuierlich.\\[0.15em]
    Ensemble, Musterfunktion: gleiche\\
    Bedeutung wie bei kont.\ Prozessen.} &
  \formelblock{CDF \& PDF (zeitabhängig)}{F_X(x,k)=P(\{\eta\mid X_\eta[k]\le x\})\\[0.15em]
    f_X(x,k)=\dfrac{\partial F_X(x,k)}{\partial x}} \\
\end{formelreihe}
\bereichende

\bereich[cDiskret]{Mittelwert, Varianz, AKF, AKV}
\begin{formelreihe}
  \formelblock{Mittelwert}{\mu_X[k]=E\{X_\eta[k]\}
    =\int_{-\infty}^{+\infty}\!x\,f_X(x,k)\,\mathrm{d}x} &
  \formelblock{Varianz}{\sigma_X^2[k]=E\{(X_\eta[k]{-}\mu_X[k])^2\}
    =\int_{-\infty}^{+\infty}\!(x{-}\mu_X)^2 f_X(x,k)\,\mathrm{d}x} &
  \formelblock{Momentanleistung}{p[k]=\varphi_{XX}[k,k]
    =E\{X_\eta^2[k]\}=\sigma_X^2[k]+\mu_X^2[k]} \\[0.3em]
  \formelblock{AKF}{\varphi_{XX}[k_1,k_2]
    =E\{X_\eta[k_1]\cdot X_\eta[k_2]\}} &
  \formelblock{AKV}{\psi_{XX}[k_1,k_2]
    =E\{(X[k_1]{-}\mu_X[k_1])(X[k_2]{-}\mu_X[k_2])\}} &
  \formelblock{Symmetrie \& Spezialfälle}{%
    \varphi_{XX}[k_1,k_2]=\varphi_{XX}[k_2,k_1]\\[0.1em]
    \varphi_{XX}[k,k]=\sigma_X^2[k]+\mu_X^2[k]\ge0\\[0.1em]
    \psi_{XX}[k,k]=\sigma_X^2[k]\ge0} \\
\end{formelreihe}
\bereichende

\bereich[cDiskret]{REZEPT · Numerische Ensemble-Auswertung ($N$ Realisierungen)}
\begin{formelreihe}
  \formelblock{Quer: Mittel-/Quadrat\-mittelwert bei festem $k$}{%
    \mu_X[k]=\dfrac1N\sum_{i=1}^{N}x_i[k]\\[0.15em]
    q_X[k]=\sqrt{\dfrac1N\sum_{i=1}^{N}x_i^2[k]}} &
  \formelblock{AKF (Ensemble) zwischen $k_1,k_2$}{%
    \varphi_{XX}[k_1,k_2]=\dfrac1N\sum_{i=1}^{N}x_i[k_1]\,x_i[k_2]} &
  \infoblock{\textbf{Quer} (Ensemble): Spalte $k$ über alle Zeilen $x_i$.\\
    \textbf{Längs} (Zeit): eine Zeile $x_i$ über alle $k$.\\[0.15em]
    \textbf{Ergodisch}: Längs-$=$Quer-Mittel. Für das Quer-Mittel braucht man
    \emph{viele} Realisierungen.} \\
\end{formelreihe}
\miniexample{$N{=}5$ Realisierungen}{%
  $q_X[7]=\sqrt{\tfrac15(0{,}57^2{+}0{,}24^2{+}0{,}42^2{+}0{,}78^2{+}0{,}68^2)}=0{,}571$;\quad
  $\varphi_{XX}[3,13]=\tfrac15\sum_i x_i[3]x_i[13]=0{,}179$.}
\bereichende

\bereich[cDiskret]{KKF, KKV — Korrelation zweier Prozesse}
\begin{formelreihe}
  \formelblock{Kreuzkorrelation (KKF)}{\varphi_{XY}[k_1,k_2]
    =E\{X_\eta[k_1]\cdot Y_\eta[k_2]\}} &
  \formelblock{Kreuzkovarianz (KKV)}{\psi_{XY}[k_1,k_2]
    =E\{(X[k_1]{-}\mu_X[k_1])(Y[k_2]{-}\mu_Y[k_2])\}} &
  \formelblock{Symmetrie}{\varphi_{XY}[k_1,k_2]=\varphi_{YX}[k_2,k_1]\\[0.1em]
    \psi_{XY}[k_1,k_2]=\psi_{YX}[k_2,k_1]} \\[0.3em]
  \formelblock{Unkorreliert}{\psi_{XY}[k_1,k_2]
    =\varphi_{XY}[k_1,k_2]-\mu_X[k_1]\mu_Y[k_2]=0} &
  \formelblock{Orthogonal}{\varphi_{XY}[k_1,k_2]=0\quad\forall\,k_1,k_2} &
  \\
\end{formelreihe}
\bereichende

\bereich[cDiskret]{Stationarität (WSS)}
\begin{formelreihe}
  \infoblock{\textbf{WSS:} statist.\ Eigenschaften unabhängig\\
    vom Abtastzeitpunkt $k$.\\[0.15em]
    $\mu_X[k]=\mu_X=\mathrm{const}$\\
    $\sigma_X^2[k]=\sigma_X^2=\mathrm{const}$\\
    $\varphi_{XX}[k_1,k_2]=\varphi_{XX}[\kappa]$\quad mit $\kappa=k_1{-}k_2$} &
  \formelblock{Stationäre AKF: $\kappa=k_1{-}k_2$}{\varphi_{XX}[\kappa]
    =E\{X[k{+}\kappa]\cdot X[k]\}\\[0.1em]
    \varphi_{XX}[0]=\sigma_X^2+\mu_X^2\quad\text{(Leistung)}\\[0.1em]
    \varphi_{XX}[-\kappa]=\varphi_{XX}[\kappa]\quad\text{(gerade)}\\[0.1em]
    |\varphi_{XX}[\kappa]|\le\varphi_{XX}[0]\quad\text{(Maximum bei }0)} &
  \formelblock{Stationäre KKF}{\varphi_{XY}[\kappa]
    =E\{X[k{+}\kappa]\cdot Y[k]\}\\[0.1em]
    \varphi_{XY}[\kappa]=\varphi_{YX}[-\kappa]} \\
\end{formelreihe}
\bereichende

\bereich[cDiskret]{Weißes Rauschen \& Ergodizität}
\begin{formelreihe}
  \infoblock{\textbf{i.i.d.\ Rauschen:} Einzelne Werte der Sequenz\\
    voneinander unabhängig und gleich verteilt\\
    (independent identically distributed).\\[0.15em]
    $\Rightarrow$ stets stationär und weiß.} &
  \formelblock{Schwach stationäres weißes Rauschen}{%
    \varphi_{XX}[\kappa]=\sigma_X^2\cdot\delta[\kappa]\\[0.15em]
    f_X(x,k)=f_X(x)\quad\text{(zeitinvariant)}} &
  \formelblock{Ergodizität}{\mu_X=\lim_{N\to\infty}\dfrac{1}{2N{+}1}
    \sum_{k=-N}^{N}x_i[k]\\[0.2em]
    \varphi_{XX}[\kappa]=\lim_{N\to\infty}\dfrac{1}{2N{+}1}
    \sum_{k=-N}^{N}x_i[k{+}\kappa]x_i[k]} \\
\end{formelreihe}
\bereichende

\bereich[cDiskret]{LDS — Diskrete Fouriertransformation (DTFT)}
\begin{formelreihe}
  \formelblock{Wiener-Khinchin (diskret)}{\Phi_{XX}(\Omega)
    =\mathrm{DTFT}\{\varphi_{XX}[\kappa]\}
    =\sum_{\kappa=-\infty}^{+\infty}\varphi_{XX}[\kappa]\,e^{-j\kappa\Omega}} &
  \formelblock{Rücktransformation}{\varphi_{XX}[\kappa]
    =\dfrac{1}{2\pi}\int_{-\pi}^{\pi}\Phi_{XX}(\Omega)\,e^{j\kappa\Omega}\,\mathrm{d}\Omega} &
  \infoblock{$\Omega=\omega T\in[-\pi,\pi]$: normierte Kreisfrequenz.\\[0.15em]
    Dimension von $\Phi_{XX}(\Omega)$: $[X]^2$\\
    (keine echte Leistungsdichte!)} \\
\end{formelreihe}
\bereichende

\bereich[cDiskret]{Diskretes LTI-System: $Y[k]=h[k]*X[k]$}
\begin{formelreihe}
  \formelblock{AKF \& KKF (Zeitbereich)}{%
    \varphi_{YX}[\kappa]=\varphi_{XX}[\kappa]*h[\kappa]\\[0.1em]
    \varphi_{XY}[\kappa]=\varphi_{XX}[\kappa]*h[-\kappa]\\[0.1em]
    \varphi_{YY}[\kappa]=\varphi_{XX}[\kappa]*h[\kappa]*h[-\kappa]\\[0.1em]
    \varphi_{hh}[\kappa]=h[\kappa]*h[-\kappa]} &
  \formelblock{LDS (Frequenzbereich)}{%
    \Phi_{YX}(\Omega)=\Phi_{XX}(\Omega)\cdot H(e^{j\Omega})\\[0.1em]
    \Phi_{XY}(\Omega)=\Phi_{XX}(\Omega)\cdot H^*(e^{j\Omega})\\[0.1em]
    \Phi_{YY}(\Omega)=\Phi_{XX}(\Omega)\cdot|H(e^{j\Omega})|^2} &
  \infoblock{Analog zum kont.\ LTI-System,\\
    aber: $H(j\omega)\to H(e^{j\Omega})$\\
    $\Omega=\omega T$,\quad$|H|^2=H\cdot H^*$\\[0.15em]
    Aus WSS-Eingang folgt\\
    WSS-Ausgang.} \\
\end{formelreihe}
\bereichende
